
\section{Introducción}

Los servicios de urgencias hospitalarias constituyen uno de los entornos clínicos con mayor exigencia cognitiva y documental del sistema sanitario. La enfermera de triage es la primera profesional que evalúa al paciente en el momento de su llegada, y en ese primer contacto, que en servicios saturados puede durar menos de cuatro minutos, debe simultáneamente registrar constantes vitales, valorar el nivel de dolor, asignar un nivel de prioridad y volcar toda esa información en el sistema de información hospitalario. Este trabajo aborda la pregunta de si la inteligencia artificial puede asumir parte de esa carga documental sin comprometer la privacidad de los datos del paciente ni requerir infraestructura en la nube.

UrgeNurse Agent propone una solución multi-agente local que integra tres fuentes de datos críticas del flujo de enfermería en urgencias: datos tabulares de triage, documentos clínicos en papel digitalizados y audio del traspaso de turno entre enfermeras. Los tres módulos se ejecutan en local, sobre CPU estándar, sin dependencia de redes externas, lo que garantiza desde el diseño el cumplimiento del Reglamento General de Protección de Datos (RGPD) y la soberanía institucional sobre los datos del paciente.


\subsection{Motivación}

La carga documental en enfermería de urgencias es un problema cuantificable y sistemáticamente documentado en la literatura científica. Un estudio de tiempo y movimiento realizado en la Vanderbilt University Medical Center con enfermeras de urgencias encontró que estas destinaban el 27\% de su turno a tareas de historia clínica electrónica, superando el tiempo dedicado a la atención directa al paciente (25\%) \parencite{bakhoum2021at}.

En el ámbito de la enfermería hospitalaria en general, una revisión en 36 hospitales identificó que el 35,3\% del tiempo de práctica enfermera se dedicaba a documentación, la actividad individual más frecuente por encima de la coordinación de cuidados o la atención directa (Hendrich et al., 2008). Una investigación separada sobre carga documental en sistemas EHR estimó que las enfermeras destinan entre un 19\% y un 35\% de su tiempo a documentación electrónica, frente al 9\% que se registraba con documentación en papel \parencite{levy2024documentation}.


Dentro de urgencias, el proceso de triage concentra riesgos específicos. MIMIC-IV-ED es la mayor base de datos p00fablica de visitas a urgencias, con más de 400.000 admisiones al servicio de urgencias del Beth Israel Deaconess Medical Center entre 2011 y 2019 (Gow et al., 2023). La literatura clínica documenta sistemáticamente la inconsistencia entre el nivel ESI registrado inicialmente y los signos vitales del paciente como uno de los principales focos de error en el proceso de triage, lo que motiva la necesidad de un sistema de validación automatizada de apoyo al criterio de la enfermera. El equipo llevará a cabo un análisis exploratorio de MIMIC-IV-ED para cuantificar este fenómeno en el dataset y reportarlo en la sección de desarrollo.
El handover de turno es el segundo punto crítico. En un servicio de urgencias activo este proceso puede cubrir simultáneamente a todos los pacientes presentes y durar entre 10 y 20 minutos. La estructura SBAR (Situation, Background, Assessment, Recommendation) es el estándar internacional recomendado por la OMS para garantizar que este traspaso sea completo y sistemático. Sin embargo, en la práctica clínica diaria el handover es frecuentemente oral, desestructurado y dependiente de la memoria del profesional. Un estudio cuasi-experimental en un servicio de urgencias en Irán encontró que la implementación del método SBAR redujo los errores clínicos totales de 102 a 25 casos en el mismo período de observación \parencite{chaharsoughi2025sbar}. En el ámbito de los handovers sin estructura SBAR, la literatura de calidad asistencial documenta pérdidas de información crítica que han sido asociadas a eventos adversos y errores de medicación \parencite{friesen2008handoffs}.

A la carga documental intrínseca se suma la heterogeneidad de los documentos clínicos que llegan al servicio desde otros centros. Analíticas de laboratorio, informes de especialistas e historiales previos en papel deben ser integrados manualmente por la enfermera en el registro activo del paciente. Cada laboratorio o clínica privada utiliza su propio formato, lo que hace imposible una integración automática sin un motor OCR capaz de manejar layouts variables. La ausencia de esta integración produce fragmentación de la información clínica que condiciona la toma de decisiones.
En este contexto, UrgeNurse Agent se posiciona no como sustituto del criterio clínico enfermero, sino como sistema de asistencia documental que procesa información existente, la estructura y la valida formalmente, manteniendo la decisión clínica bajo supervisión profesional. Desde la perspectiva del AI Act europeo, esta arquitectura sitúa el sistema en la categoría de soporte a la decisión con supervisión humana obligatoria, con requisitos de transparencia y trazabilidad, pero sin necesidad de validación clínica formal propia de sistemas de diagnóstico autónomo.


\subsection{Planteamiento del trabajo}

El problema identificado tiene tres dimensiones documentales distintas pero interrelacionadas: la validación del triage inicial, la integración de documentación clínica externa y la estructuración del handover de turno. Ninguna herramienta existente en el ecosistema open source aborda las tres de forma integrada, ejecutándose en local y sobre datos reales de urgencias. Las soluciones actuales tratan cada modalidad por separado, y las que usan LLMs requieren el envío de datos a servicios en la nube, lo que las hace incompatibles con los requisitos de privacidad del entorno sanitario europeo.
UrgeNurse Agent propone un pipeline multi-agente que procesa las tres fuentes de datos más críticas del flujo de enfermería en urgencias —CSV de triage, PDF de analíticas y audio de handover— para generar tres outputs estandarizados: triage ESI validado con alertas de inconsistencia, nota nursing estructurada con nivel de alerta analítica, y resumen SBAR de handover, ejecutando todos los modelos en local sin dependencia de red ni de servicios externos.

La solución técnica se articula sobre tres decisiones de diseño fundamentales que se justifican en detalle en la sección de desarrollo. Primera, la orquestación mediante LangGraph permite que los tres agentes compartan un estado tipado (PacienteState) y se ejecuten en paralelo, con soporte nativo para condicionales y re intentos. Segunda, el uso de un único modelo de lenguaje local (Mistral-7B vía Ollama) para todas las tareas generativas garantiza privacidad sin sacrificar rendimiento en CPU estándar. Tercera, la evaluación de cada módulo se realiza con métricas cuantitativas específicas (kappa de Cohen para triage, F1-score para OCR, WER para ASR) antes de evaluar el sistema integrado, siguiendo el principio de que una evaluación agregada sin contexto no tiene valor académico.

\subsection{Estructura del trabajo}

La sección 2 sitúa el problema en su contexto clínico y tecnológico mediante análisis de la carga documental en enfermería de urgencias extraído de la literatura científica disponible. Se realiza una revisión crítica de las soluciones existentes en tres ejes: sistemas de triage asistido por inteligencia artificial, motores de OCR aplicados a documentación clínica heterogénea, y modelos de reconocimiento automático del habla en entornos médicos. La sección concluye identificando la brecha de integración que ninguna solución actual cubre: la orquestación local de las tres modalidades en un pipeline único orientado a la enfermería de urgencias.

La sección 3 define el objetivo general y los seis objetivos específicos con sus métricas de éxito cuantitativas. Describe la metodología de trabajo adoptada para el desarrollo del proyecto, incluyendo el proceso de acceso a datasets, las decisiones de diseño arquitectónico y el plan de evaluación de cada módulo.
La sección 4 constituye el núcleo técnico del trabajo. Se desarrollan los tres módulos de UrgeNurse Agent de forma independiente —Agente Triage ESI, Agente Documental OCR y Agente ASR+SBAR— describiendo para cada uno las decisiones de implementación, la comparativa de alternativas y los resultados de evaluación. La sección concluye con la integración de los tres agentes en la arquitectura LangGraph y la evaluación del sistema completo sobre escenarios clínicos end-to-end.
La sección 5 recoge las conclusiones del trabajo en relación con los objetivos planteados, discute las limitaciones del sistema y propone las líneas de trabajo futuro más relevantes para evolucionar UrgeNurse Agent hacia un entorno de producción hospitalario real.
